\documentclass[12pt,a4paper]{article}

% ─── Paquetes ───────────────────────────────────────────────────────────────
\usepackage[utf8]{inputenc}
\usepackage[T1]{fontenc}
\usepackage[spanish]{babel}
\usepackage{amsmath,amssymb}
\usepackage{booktabs}
\usepackage{array}
\usepackage{geometry}
\usepackage{xcolor}
\usepackage{titlesec}
\usepackage{enumitem}
\usepackage{fancyhdr}
\usepackage{mdframed}
\usepackage{multicol}
\usepackage{tasks}

% ─── Geometría ──────────────────────────────────────────────────────────────
\geometry{left=2.5cm, right=2.5cm, top=3cm, bottom=3cm}

% ─── Colores ────────────────────────────────────────────────────────────────
\definecolor{primario}{RGB}{0,90,160}
\definecolor{secundario}{RGB}{220,235,250}
\definecolor{acento}{RGB}{200,50,50}
\definecolor{grisclaro}{RGB}{245,245,245}

% ─── Encabezado / pie ────────────────────────────────────────────────────────
\pagestyle{fancy}
\fancyhf{}
\fancyhead[L]{\small\textcolor{primario}{\textbf{Ejercicios — Conectivos Lógicos}}}
\fancyhead[R]{\small\textcolor{primario}{Cuarto Grado de Secundaria}}
\fancyfoot[C]{\thepage}
\renewcommand{\headrulewidth}{0.6pt}

% ─── Estilos de título ───────────────────────────────────────────────────────
\titleformat{\section}{\large\bfseries\color{primario}}{}{0pt}{}[\titlerule]
\titleformat{\subsection}{\normalsize\bfseries\color{acento}}{}{0pt}{}

% ─── Entorno para recuadro de ejercicio ──────────────────────────────────────
\newmdenv[
  backgroundcolor=grisclaro,
  linecolor=primario,
  linewidth=1pt,
  leftmargin=0pt, rightmargin=0pt,
  innerleftmargin=10pt, innerrightmargin=10pt,
  innertopmargin=6pt, innerbottommargin=6pt
]{ejerciciobox}

% ─── Espacio para respuesta ──────────────────────────────────────────────────
\newcommand{\linea}{\underline{\hspace{10cm}}}
\newcommand{\espacioresp}[1][3cm]{\vspace{#1}}

% ─── Documento ───────────────────────────────────────────────────────────────
\begin{document}

% Portada
\begin{center}
  {\LARGE\bfseries\color{primario} CONECTIVOS LÓGICOS Y TABLA DE VERDAD}\\[4pt]
  {\large Ejercicios — Cuarto Grado de Secundaria}\\[2pt]
  \textcolor{gray}{\rule{0.6\linewidth}{0.6pt}}\\[6pt]
  Nombre: \underline{\hspace{7cm}} \quad Fecha: \underline{\hspace{3cm}}
\end{center}

\vspace{6pt}

% ════════════════════════════════════════════════════════════════════════════
\section{Parte I — Nivel Integral}
% ════════════════════════════════════════════════════════════════════════════

% Ejercicio 1
\begin{ejerciciobox}
\textbf{Ejercicio 1.}
Determina la \textbf{matriz principal} de la siguiente proposición compuesta:
\[
(p \wedge q) \vee q
\]
\end{ejerciciobox}

\espacioresp[5cm]

% Ejercicio 2
\begin{ejerciciobox}
\textbf{Ejercicio 2.}
Determina los valores de verdad de $r$ y $p$ si se sabe que la siguiente
proposición es \textbf{falsa}:
\[
\sim p \vee r
\]
\end{ejerciciobox}

\espacioresp[3.5cm]

% Ejercicio 3
\begin{ejerciciobox}
\textbf{Ejercicio 3.}
Señala la \textbf{proposición compuesta} entre las siguientes opciones:

\begin{enumerate}[label=\alph*)]
  \item Agripino y Cesarina son hermanos.
  \item Los Heraldos Negros es una obra de César Vallejo.
  \item Joseph-Nicéphore tomó la primera fotografía en blanco y negro.
  \item Carlos y Richard van juntos al cine.
  \item Daniel es profesor y Rosa es escritora.
\end{enumerate}
\end{ejerciciobox}

\espacioresp[1.5cm]

% ════════════════════════════════════════════════════════════════════════════
\section{Parte II — Simbolización Lógica}
% ════════════════════════════════════════════════════════════════════════════

% Ejercicio 4
\begin{ejerciciobox}
\textbf{Ejercicio 4.}
Simboliza mediante conectivos lógicos la siguiente proposición:

\begin{center}
  \textit{``Si Daniel y Agripina juegan fútbol, Margarito será el árbitro.''}
\end{center}

Define las variables:
\begin{itemize}
  \item $p$: Daniel juega fútbol.
  \item $q$: Agripina juega fútbol.
  \item $r$: Margarito será el árbitro.
\end{itemize}

Simbolización: \linea
\end{ejerciciobox}

\espacioresp[2cm]

% Ejercicio 5
\begin{ejerciciobox}
\textbf{Ejercicio 5.}
Simboliza mediante conectivos lógicos la siguiente proposición:

\begin{center}
  \textit{``Si tomas jugo de naranja o fresa, entonces estarás lleno.''}
\end{center}

Define tus propias variables y escribe la simbolización:

\vspace{0.5cm}
Variables: \linea

Simbolización: \linea
\end{ejerciciobox}

\espacioresp[2cm]

% ════════════════════════════════════════════════════════════════════════════
\section{Parte III — Matrices Principales}
% ════════════════════════════════════════════════════════════════════════════

% Ejercicio 6
\begin{ejerciciobox}
\textbf{Ejercicio 6.}
Determina la \textbf{matriz principal} de la siguiente proposición compuesta:
\[
(p \;\Delta\; q) \leftrightarrow \sim r
\]
Completa la tabla de verdad ($2^3 = 8$ filas):

\vspace{0.4cm}
\begin{center}
\begin{tabular}{ccc|c|c|c|c}
\toprule
$p$ & $q$ & $r$ & $p \;\Delta\; q$ & $\sim r$ & $(p \;\Delta\; q) \leftrightarrow \sim r$ \\
\midrule
V & V & V & & & \\[8pt]
V & V & F & & & \\[8pt]
V & F & V & & & \\[8pt]
V & F & F & & & \\[8pt]
F & V & V & & & \\[8pt]
F & V & F & & & \\[8pt]
F & F & V & & & \\[8pt]
F & F & F & & & \\[8pt]
\bottomrule
\end{tabular}
\end{center}
\textbf{Matriz principal:} \linea
\end{ejerciciobox}

\newpage

% Ejercicio 7
\begin{ejerciciobox}
\textbf{Ejercicio 7.}
Si la proposición compuesta:
\[
\bigl[(p \rightarrow q) \vee (q \vee \sim r)\bigr]
\]
es \textbf{falsa}, determina los valores de verdad de $p$, $q$ y $r$.
\end{ejerciciobox}

\espacioresp[5cm]

% ════════════════════════════════════════════════════════════════════════════
\section{Parte IV — Nivel UNMSM / UNI}
% ════════════════════════════════════════════════════════════════════════════

% Ejercicio 8
\begin{ejerciciobox}
\textbf{Ejercicio 8. \small(Nivel UNMSM)}
Si la siguiente proposición es \textbf{falsa}:
\[
(\sim p \wedge q) \rightarrow \bigl[(p \vee r) \vee t\bigr]
\]
determina el valor de verdad de:
\begin{enumerate}[label=\Roman*.]
  \item $\sim(\sim p \vee \sim q) \rightarrow (r \vee \sim t)$
  \item $(\sim p \rightarrow t) \rightarrow (\sim q \rightarrow r)$
\end{enumerate}
\end{ejerciciobox}

\espacioresp[6cm]

% Ejercicio 9
\begin{ejerciciobox}
\textbf{Ejercicio 9.}
Si la proposición:
\[
(p \rightarrow \sim q) \vee (\sim r \rightarrow s)
\]
es \textbf{falsa}, determina los valores de verdad de las siguientes proposiciones:
\begin{enumerate}[label=\Roman*.]
  \item $(\sim p \wedge \sim q) \vee \sim p$
  \item $(p \rightarrow q) \rightarrow r$
\end{enumerate}
\end{ejerciciobox}

\espacioresp[6cm]

% Ejercicio 10
\begin{ejerciciobox}
\textbf{Ejercicio 10.}
Determina si la siguiente proposición es \textbf{tautología}, \textbf{contradicción} o
\textbf{contingente}:
\[
\bigl[(\sim p \wedge q) \rightarrow r\bigr] \leftrightarrow \bigl[(p \wedge q) \;\Delta\; \sim r\bigr]
\]
\end{ejerciciobox}

\espacioresp[6cm]

\newpage

% Ejercicio 11
\begin{ejerciciobox}
\textbf{Ejercicio 11.}
Determina el valor de verdad de las siguientes proposiciones:
\begin{enumerate}[label=\alph*)]
  \item $(3 + 5 = 9) \wedge (5 \times 2 = 10)$
  \item $\left(\dfrac{13}{5} + 1 = \dfrac{18}{5}\right) \rightarrow (3^2 = 5)$
  \item $(2^3 = 8) \;\Delta\; (\sqrt{16} = -4)$
  \item $(-13 < 8) \leftrightarrow (8 + 1 > 9)$
\end{enumerate}
\end{ejerciciobox}

\espacioresp[5cm]

% Ejercicio 12
\begin{ejerciciobox}
\textbf{Ejercicio 12. \small(UNI 2013-I)}
Dada la proposición:
\[
\sim\bigl[(r \vee q) \rightarrow (r \rightarrow p)\bigr] \equiv V
\]
donde $q$ es una proposición \textbf{falsa}. Determina el valor de verdad de:
\begin{enumerate}[label=\Roman*.]
  \item $r \rightarrow (\sim p \vee \sim q)$
  \item $\bigl[r \leftrightarrow (p \vee q)\bigr] \leftrightarrow (q \wedge \sim p)$
  \item $(r \vee \sim p) \wedge (q \vee p)$
\end{enumerate}
\end{ejerciciobox}

\espacioresp[6cm]

% Ejercicio 13
\begin{ejerciciobox}
\textbf{Ejercicio 13. \small(UNI 2012-II)}
Si la proposición:
\[
\bigl[(\sim p \vee q) \rightarrow (q \leftrightarrow r)\bigr] \vee (q \wedge s)
\]
es \textbf{falsa} y $p$ es una proposición \textbf{verdadera}, determina los valores de
verdad de $q$, $r$ y $s$ (en ese orden).
\end{ejerciciobox}

\espacioresp[5cm]

% Ejercicio 14
\begin{ejerciciobox}
\textbf{Ejercicio 14.}
Clasifica las siguientes proposiciones como tautología (T), contradicción (F)
o contingencia (C):
\begin{enumerate}[label=\Roman*.]
  \item $(p \rightarrow q) \rightarrow \sim q$
  \item $(\sim q \vee q) \;\Delta\; \bigl[p \;\Delta\; (p \vee q)\bigr]$
  \item $(q \;\Delta\; \sim p) \leftrightarrow (p \;\Delta\; q)$
\end{enumerate}
\end{ejerciciobox}

\espacioresp[5cm]

\end{document}
